\documentclass{article}


% if you need to pass options to natbib, use, e.g.:
%     \PassOptionsToPackage{numbers, compress}{natbib}
% before loading neurips_2024


% ready for submission
\usepackage[final]{neurips_2024}


% to compile a preprint version, e.g., for submission to arXiv, add add the
% [preprint] option:
%     \usepackage[preprint]{neurips_2024}


% to compile a camera-ready version, add the [final] option, e.g.:
%     \usepackage[final]{neurips_2024}


% to avoid loading the natbib package, add option nonatbib:
%    \usepackage[nonatbib]{neurips_2024}


\usepackage[utf8]{inputenc} % allow utf-8 input
\usepackage[T1]{fontenc}    % use 8-bit T1 fonts
\usepackage{hyperref}       % hyperlinks
\usepackage{url}            % simple URL typesetting
\usepackage{booktabs}       % professional-quality tables
\usepackage{amsfonts}       % blackboard math symbols
\usepackage{nicefrac}       % compact symbols for 1/2, etc.
\usepackage{microtype}      % microtypography
\usepackage{xcolor}         % colors


\title{Comparing Neural Network Architectures for sEMG-to-Keystroke Prediction}


\author{%
  Student A\thanks{} \\
  Department of Electrical and Computer Engineering\\
  University of California, Los Angeles\\
  Los Angeles, CA 90095 \\
  \texttt{wesleygwn@ucla.edu, ahjcuniv@ucla.edu} \\
}


\begin{document}


\maketitle


\begin{abstract}
  We compare four neural network architectures: TDS Conv, GRU, LSTM, and CNN+RNN for predicting keystrokes from forearm sEMG signals using the emg2qwerty dataset and CTC loss. All models were trained on single-user 16-channel recordings across two hardware configurations. On Config. A (2$\times$RTX 3060), LSTM CTC achieves the best test CER of 15.99\% using 2$\times$ temporal downsampling (2 kHz $\to$ 1 kHz) across 11 unique configurations including electrode ablations. On Config. B (RTX 4080) at full 2 kHz sampling rate with matching configuration (batch=32, epochs=150), LSTM CTC also performs best at 18.59\%. Cross-platform evaluation confirms reproducibility: TDS Conv yields 22.28\% (Config. A) vs 21.79\% (Config. B), within 3\% relative difference. The best LSTM configuration uses full 16-channel input (1056 features), with performance degrading substantially when reducing electrode channels.
\end{abstract}


\section{Introduction}

Surface electromyography (sEMG) records electrical activity from muscles through skin-placed electrodes, offering a non-invasive way to capture muscle signals for human-computer interaction. The forearm is a natural target for sEMG-based typing detection because finger and wrist movements activate distinguishable muscle groups.

Sivakumar et~al.\ \citep{sivakumar2024emg2qwertylargedatasetbaselines} introduced emg2qwerty, a dataset of 1,135 recording sessions spanning 108 users and 346 hours of 16-channel sEMG recordings with ground-truth keystrokes. The authors provided a Time-Delay Separable Convolution (TDS Conv) baseline trained with Connectionist Temporal Classification (CTC) loss, which handles alignment by marginalizing over all frame-to-character mappings.

The choice of architecture for sequence-to-sequence tasks remains open. Convolutional approaches like TDS Conv capture local temporal patterns efficiently, while recurrent architectures (LSTMs, GRUs) model longer-range dependencies, each with different inductive biases.

This work compares four neural network architectures for sEMG-to-keystroke prediction: TDS Conv CTC, GRU CTC, LSTM CTC, and CNN+RNN CTC. All models use CTC loss and are evaluated on single-user data. We measure performance using character error rate (CER) and error type breakdown. To assess reproducibility and explore architectural variations, we run experiments on two hardware configurations: Config. A (2$\times$RTX 3060, 11 unique model configurations including electrode ablations) and Config. B (RTX 4080, 4 architectures).

\section{Methods}

\subsection{Dataset}

We use the emg2qwerty dataset \citep{sivakumar2024emg2qwertylargedatasetbaselines}, which contains surface electromyography (sEMG) recordings from users touch typing on a QWERTY keyboard. The dataset comprises 1,135 recording sessions spanning 108 users and 346 hours of data. Each session includes 16-channel sEMG recordings from both left and right forearms, sampled at 2 kHz, along with ground-truth keystrokes captured via keylogger.

For this work, we focus on single-user training to establish a controlled comparison between architectures. We use the default data splits provided by the dataset: training sessions for model optimization, validation sessions for hyperparameter tuning and early stopping, and test sessions for final evaluation. The sEMG signals are recorded using two Delsys Trigno wireless electrodes placed on each forearm, yielding 16 channels per arm (8 per electrode). The recordings capture muscle activations from flexor and extensor muscle groups involved in finger and wrist movements during typing.

\subsection{Preprocessing}

The raw sEMG signals undergo several preprocessing steps before being fed to the neural networks. First, the signals are converted from numpy structured arrays to PyTorch tensors, stacking the left and right forearm channels along a band dimension to produce tensors of shape $(T, 2, 16)$, where $T$ is the temporal dimension, 2 represents the two forearm bands, and 16 is the number of electrode channels.

We compute log-scaled spectrograms from the time-domain signals using Short-Time Fourier Transform (STFT) with an FFT size of 64 (yielding 33 frequency bins) and a hop length of 16 samples. This produces time-frequency representations of shape $(T, 2, 16, 33)$, where the spectrogram is computed independently for each channel. The log transformation ($\log_{10}(x + 10^{-6})$) stabilizes training by compressing the dynamic range of the spectral features.

\paragraph{Configuration-Specific Preprocessing.} Most models use the full 2 kHz sampling rate without downsampling. As an experimental variation, two LSTM models on Config. A explore temporal downsampling: one uses $2\times$ downsampling (2 kHz $\to$ 1 kHz) and another uses $4\times$ downsampling (2 kHz $\to$ 500 Hz). All other Config. A models and all Config. B models use the full 2 kHz sampling rate.

During training, we apply data augmentation techniques to improve generalization: random band rotation (shifting electrode channels by $\{-1, 0, 1\}$ positions) to simulate electrode placement variability, temporal alignment jitter (up to 20 samples) between left and right forearm signals, and SpecAugment with time and frequency masking. The models receive 4-second windows (8000 samples at 2 kHz) with 900 ms of past context and 100 ms of future context padding.

\subsection{Model Architectures}

We compare four neural network architectures, all trained with Connectionist Temporal Classification (CTC) loss. Each architecture shares a common frontend for feature extraction but differs in its temporal modeling approach.

\paragraph{TDS Conv CTC.} The baseline architecture uses Time-Delay Separable (TDS) convolutions for temporal modeling. The frontend consists of spectrogram normalization and a multi-band rotation-invariant MLP that processes each forearm band independently before flattening. The TDS encoder comprises 4 convolutional blocks with 24 channels each and a kernel width of 32, yielding a total temporal receptive field of 125 samples. This architecture efficiently captures local temporal patterns while maintaining parameter efficiency through depthwise separable convolutions.

\paragraph{GRU CTC.} This architecture replaces the TDS convolutional encoder with a bidirectional Gated Recurrent Unit (GRU). After the shared frontend (MLP with hidden features), a bidirectional GRU models temporal dependencies. The bidirectional design allows the model to incorporate both past and future context when making predictions at each timestep. Dropout is applied between recurrent layers.

\paragraph{LSTM CTC.} Similar to the GRU variant, this model uses recurrent layers but with Long Short-Term Memory (LSTM) cells. The architecture employs a bidirectional LSTM with dropout. The LSTM's more expressive gating mechanism (input, forget, and output gates) provides greater capacity for modeling complex temporal patterns compared to GRUs.

\paragraph{CNN+RNN CTC.} This hybrid architecture combines convolutional feature extraction with recurrent modeling. After the frontend MLP, temporal convolutions extract local features. A bidirectional GRU then models longer-range dependencies. This design aims to leverage the complementary strengths of both architectural paradigms.

All architectures conclude with a linear layer mapping to the character set (80 classes including letters, digits, punctuation, and special tokens) followed by log-softmax for CTC training.

\subsection{Training Setup}

All models are trained using Connectionist Temporal Classification (CTC) loss, which enables sequence-to-sequence learning without explicit alignment between input frames and output characters. CTC marginalizes over all valid alignments by introducing a blank token and using a dynamic programming algorithm to compute the loss efficiently.

We use the Adam optimizer with a learning rate of $10^{-3}$. The learning rate schedule combines linear warmup over 10 epochs followed by cosine annealing decay to a minimum learning rate of $10^{-6}$. This schedule helps stabilize early training while allowing fine-grained adaptation in later epochs.

\subsection{Experimental Configurations}

We conduct experiments on two hardware configurations to assess reproducibility:

\textbf{Configuration A:} Batch size 32, maximum 150 epochs, trained on 2$\times$NVIDIA RTX 3060 GPUs (12GB each) via local Jupyter server.

\textbf{Configuration B:} Batch size 32, maximum 150 epochs, trained on a single NVIDIA RTX 4080 GPU.

For evaluation, we use character error rate (CER) computed via the Levenshtein distance between predicted and ground-truth character sequences. We additionally report insertion error rate (IER), deletion error rate (DER), and substitution error rate (SER) to provide fine-grained analysis of model behavior. Greedy decoding is used at test time for all experiments.

\section{Results}

\subsection{Overall Results}

Table~\ref{tab:main_results} presents all configurations trained on both platforms with batch size 32 and maximum 150 epochs. The best overall performance is achieved by LSTM on Config. A with $2\times$ temporal downsampling (2 kHz $\to$ 1 kHz), achieving a test CER of 15.99\%.

\begin{table}[h]
\centering
\caption{Overall results across all configurations (batch=32, epochs=150). Best test CER shown in bold.}
\label{tab:main_results}
\begin{tabular}{lllccccc}
\toprule
\textbf{Architecture} & \textbf{Config} & \textbf{Downsample} & \textbf{Channels} & \textbf{Test CER} & \textbf{IER} & \textbf{DER} & \textbf{SER} \\
\midrule
LSTM & A & 2$\times$ & 16 & \textbf{15.99\%} & 2.81\% & 1.45\% & 11.74\% \\
LSTM & A & --- & 16 & 17.05\% & 2.83\% & 1.84\% & 12.38\% \\
LSTM & B & --- & 16 & 18.59\% & 2.25\% & 2.64\% & 13.70\% \\
CNN+RNN & A & --- & 16 & 19.04\% & 3.89\% & 2.29\% & 12.86\% \\
LSTM & A & 4$\times$ & 16 & 19.88\% & 7.48\% & 0.76\% & 11.65\% \\
LSTM & A & --- & 16 & 20.36\% & 3.59\% & 2.59\% & 14.18\% \\
TDS Conv & B & --- & 16 & 21.79\% & 5.12\% & 2.16\% & 14.50\% \\
TDS Conv & A & --- & 16 & 22.28\% & 5.36\% & 1.75\% & 15.17\% \\
LSTM & A & --- & 16 & 22.52\% & 4.65\% & 2.14\% & 15.73\% \\
CNN+LSTM & A & --- & 16 & 22.71\% & 3.80\% & 2.27\% & 16.64\% \\
LSTM & A & --- & 8 & 23.08\% & 4.32\% & 1.69\% & 17.07\% \\
GRU & B & --- & 16 & 27.04\% & 5.92\% & 2.94\% & 18.18\% \\
GRU & A & --- & 16 & 29.00\% & 4.99\% & 3.26\% & 20.75\% \\
LSTM & A & --- & 4 & 35.34\% & 8.23\% & 1.99\% & 25.11\% \\
LSTM & A & --- & 1 & 54.61\% & 25.09\% & 0.86\% & 28.66\% \\
\bottomrule
\end{tabular}
\end{table}

Multiple LSTM runs with the same configuration show substantial variance (15.99\% to 22.52\%), indicating sensitivity to random initialization. The electrode channel ablation study confirms that full 16-channel input is critical for optimal performance.

\subsection{Cross-Platform Comparison}

Table~\ref{tab:cross_platform} compares LSTM CTC performance across both configurations using matching settings (batch=32, epochs=150). Note that most Config. A models use the full 2 kHz sampling rate, with only one run using $2\times$ downsampling (2 kHz $\to$ 1 kHz) and another using $4\times$ downsampling (2 kHz $\to$ 500 Hz). Config. B uses the full 2 kHz sampling rate without downsampling.

\begin{table}[h]
\centering
\caption{Cross-platform LSTM CTC comparison (batch=32, epochs=150).}
\label{tab:cross_platform}
\begin{tabular}{llcccc}
\toprule
\textbf{Architecture} & \textbf{Config} & \textbf{Test CER} & \textbf{Val CER} & \textbf{IER} & \textbf{Hardware} \\
\midrule
LSTM & A & \textbf{17.05\%} & 17.86\% & 2.83\% & 2$\times$RTX 3060 \\
LSTM & B & 18.59\% & 16.97\% & 2.25\% & RTX 4080 \\
\bottomrule
\end{tabular}
\end{table}

The Config. A LSTM (17.05\%) outperforms the Config. B run (18.59\%) by 1.54 percentage points. This gap may be attributable to random initialization variance or minor implementation differences. Both configurations confirm LSTM as the best-performing architecture.

\subsection{Electrode Channel Ablation}

All ablation experiments use the following hyperparameters: no downsampling (full 2 kHz sampling rate), batch size 32, maximum 150 epochs, learning rate $10^{-3}$, Adam optimizer with linear warmup over 10 epochs followed by cosine annealing decay.

Table~\ref{tab:ablation} shows the impact of reducing electrode channels on LSTM performance.

\begin{table}[h]
\centering
\caption{LSTM electrode channel ablation (Config. A). Performance degrades substantially with fewer channels.}
\label{tab:ablation}
\begin{tabular}{lclcccc}
\toprule
\textbf{Config} & \textbf{Channels} & \textbf{Features} & \textbf{Test CER} & \textbf{IER} & \textbf{SER} \\
\midrule
A & 16 & 1056 & 15.99\% & 2.81\% & 11.74\% \\
A & 8 & 528 & 23.08\% & 4.32\% & 17.07\% \\
A & 4 & 264 & 35.34\% & 8.23\% & 25.11\% \\
A & 1 & 66 & 54.61\% & 25.09\% & 28.66\% \\
\bottomrule
\end{tabular}
\end{table}

The ablation study confirms that spatial information from multiple electrodes is essential. Reducing from 16 to 8 channels increases CER by 7.09 percentage points. Further reductions show catastrophic degradation: 4 channels yields 35.34\% CER and single channel yields 54.61\% CER.

\section{Discussion}

\subsection{LSTM Achieves Best Overall Performance}

The LSTM CTC model with full 16-channel input achieves the best overall test CER of 15.99\% on Config. A, outperforming all other configurations by a substantial margin. This result demonstrates that LSTM's more expressive gating mechanism (with separate input, forget, and output gates plus explicit cell state) effectively captures the temporal dependencies in sEMG signals during typing.

The multiple runs of the same LSTM configuration (15.99\%, 19.88\%, 20.36\%, 22.52\%) reveal significant variance from random initialization, with a range of 6.53 percentage points. The best run (15.99\%) substantially outperforms the others, suggesting that careful initialization or multiple restarts may be necessary to achieve optimal LSTM performance.

\subsection{GRU vs LSTM Analysis}

The comparison between GRU and LSTM reveals significant performance differences. On Config. A (most models at full 2 kHz, with one run using $2\times$ downsampling), LSTM achieves 15.99\% while GRU performs poorly at 29.00\%. On Config. B (without downsampling) with matching configuration (batch=32, epochs=150), LSTM achieves 18.59\% while GRU performs worse at 27.04\%.

This discrepancy suggests that GRU may require more careful hyperparameter tuning or different initialization strategies to achieve competitive performance. LSTM consistently outperforms GRU across both configurations when using the same configuration. Note that performance comparisons across configurations are confounded by the preprocessing difference for the best Config. A run.

\subsection{Cross-Platform Variance}

The TDS Conv baseline shows a 0.61 percentage point difference between Config. A (22.28\%) and Config. B (22.89\%) configurations. This variance falls within expected experimental bounds given multiple sources of non-determinism: different GPU architectures (Ampere vs Ada Lovelace), random initialization effects. Both use full 2 kHz sampling rate without downsampling.

The comparable results confirm that TDS Conv performance is reproducible across hardware platforms. The GRU discrepancy (29.00\% on Config. A vs 27.04\% on Config. B with matching configuration) suggests that some architectures may be more sensitive to hyperparameter choices and initialization. Note that performance comparisons across configurations are confounded by the preprocessing difference for the best Config. A run.

\subsection{Electrode Channel Importance}

The ablation study demonstrates that spatial information from multiple electrodes is essential for accurate keystroke prediction. The near-linear degradation in performance as channels are reduced suggests that each electrode provides complementary information about muscle activations.

The dramatic increase in insertion error rate (from 3.03\% to 25.09\% when reducing from 16 to 1 channel) indicates that without sufficient spatial resolution, the model becomes highly uncertain and prone to spurious character predictions. This has practical implications for hardware design: fewer electrodes reduce cost and complexity but significantly impact accuracy.

\subsection{Limitations}

This work has several limitations. First, the cross-configuration comparison is confounded by multiple simultaneous differences (hardware, random seeds, and preprocessing for specific models), making it difficult to isolate the effect of any single factor.

Second, most Config. A models use the full 2 kHz sampling rate, but two LSTM variants explore temporal downsampling (2$\times$ and 4$\times$ respectively). Config. B uses the full 2 kHz sampling rate without downsampling. This makes direct cross-platform performance comparison confounded when comparing the downsampled Config. A runs to Config. B models.

Third, we did not conduct systematic hyperparameter tuning. The discrepancy between GRU performance on Config. A (29.00\%) vs Config. B (27.04\%) suggests that hyperparameter choices significantly impact results. A systematic hyperparameter search could improve absolute performance for each architecture.

Fourth, all experiments were conducted on single-user data. The emg2qwerty dataset supports multi-user training, and the generalization of these architectural comparisons to cross-user scenarios is unknown.

Fifth, we used greedy decoding only. Beam search decoding with language model rescoring could improve absolute performance but was not evaluated.

\subsection{Practical Implications}

The experimental results suggest different architectures may be preferred depending on application constraints. For accuracy-critical applications where training time is not a primary concern, LSTM CTC with full 16-channel input provides the best test performance (15.99\% on Config. A using $2\times$ downsampling) when properly initialized. However, the high variance across runs (6.53 percentage points) suggests multiple training restarts may be necessary.

For real-time or resource-constrained applications, TDS Conv offers a good trade-off between accuracy (22.28\% on Config. A at full 2 kHz, 21.79\% on Config. B at full 2 kHz) and inference speed. Its purely convolutional architecture enables efficient GPU utilization and lower latency.

The GRU performs worse than LSTM across all configurations (27.04-29.00\% CER) and may be more sensitive to hyperparameter choices. Note that performance comparisons between Config. A (with one run using $2\times$ downsampling) and Config. B (without downsampling) are confounded by this preprocessing difference.

\begin{ack}
This work was completed as part of the ECE C147/247A course project at UCLA.
\end{ack}

\section*{References}


{
\small

\bibliographystyle{plain}
\begin{thebibliography}{99}

\bibitem[Cho et~al.(2014)]{cho2014learning}
Cho, K., Van~M{\"e}rri{\"e}nboer, B., Gulcehre, C., Bahdanau, D., Bougares, F., Schwenk, H., \& Bengio, Y. (2014).
\newblock Learning phrase representations using rnn encoder-decoder for statistical machine translation.
\newblock \emph{arXiv preprint arXiv:1406.1078}.

\bibitem[Graves et~al.(2006)]{graves2006connectionist}
Graves, A., Fern{\'a}ndez, S., Gomez, F., \& Schmidhuber, J. (2006).
\newblock Connectionist temporal classification: Labelling unsegmented sequence data with recurrent neural networks.
\newblock In \emph{Proceedings of the 23rd International Conference on Machine Learning} (pp. 369--376).

\bibitem[Hochreiter \& Schmidhuber(1997)]{hochreiter1997long}
Hochreiter, S., \& Schmidhuber, J. (1997).
\newblock Long short-term memory.
\newblock \emph{Neural Computation}, 9(8), 1735--1780.

\bibitem[Sivakumar et~al.(2024)]{sivakumar2024emg2qwertylargedatasetbaselines}
Sivakumar, V., Seely, J., Du, A., Bittner, S.~R., Berenzweig, A., Bolarinwa, A., Gramfort, A., \& Mandel, M.~I. (2024).
\newblock emg2qwerty: A large dataset with baselines for touch typing using surface electromyography.
\newblock \emph{arXiv preprint arXiv:2410.20081}.

\end{thebibliography}
}


\end{document}
